\section{The Composition That Passes}
\label{sec:ch15-the-composition-that-passes}

\index{resolvability!the graph that composes and cannot be served|(}%
Section~\ref{sec:ch15-the-route-it-named-is-one-of-six} took Speakers out of
\code{graph.yaml} and the composer still refused. Put Speakers back and take
Search out instead, which leaves the three services
chapter~\ref{ch:second-third-subgraph} finished with, before Search existed.
Same broken Ratings document, one subgraph entry different from a graph that
has just produced three errors. Compose it.

\begin{minted}[fontsize=\footnotesize,breakanywhere]{text}
Router execution config successfully written to ".../graph/router.json".
\end{minted}

\index{resolvability!exit zero on an unservable graph}%
It composes. Exit zero, nothing printed, and no mention of the three fields
that were unresolvable in the graph this one differs from by a single
subgraph.

\index{router execution config!carries the contradiction openly}%
The config it wrote is not hiding anything. Chapter~\ref{ch:composition} opened
one of these and read the datasource entries; the Ratings entry now holds both
halves of the contradiction, in one object:

\begin{minted}[fontsize=\footnotesize,breakanywhere]{json}
"keys": [{"typeName": "Session", "selectionSet": "id",
          "disableEntityResolver": true}],
"rootNodes": [{"typeName": "Query", "fieldNames": ["ratingCount"]},
              {"typeName": "Session",
               "fieldNames": ["averageScore", "ratingCount",
                              "feedbackUrl", "id"],
               "externalFieldNames": ["title"]}]
\end{minted}

\index{router execution config!three fields and no door}%
Three fields the router is told this service resolves, and beside them the flag
chapter~\ref{ch:second-third-subgraph} introduced, saying the router may not
send it an entity. The composer wrote both, in one file, in one pass, and
called the graph good.

\subsection{What Buys the Silence}

\index{node@\code{node(id:)}!the field that changes the answer}%
One document decides whether these three subgraphs pass or fail, and it is not
the broken one. It is Speakers, which declares no \code{Session} anywhere in
it. Change it one piece at a time, holding the other two fixed, and the check
turns on and off:

\begin{center}
\begin{tabular}{l l}
\toprule
The Speakers document declares & Composition \tabularnewline
\midrule
\code{node} and \code{nodes}, \code{Speaker implements Node} & passes
  \tabularnewline
\addlinespace
neither, \code{implements Node} kept & fails \tabularnewline
\addlinespace
neither, \code{implements Node} removed & fails \tabularnewline
\addlinespace
both, \code{implements Node} removed & passes \tabularnewline
\addlinespace
\code{node} alone & passes \tabularnewline
\addlinespace
\code{nodes} alone & fails \tabularnewline
\bottomrule
\end{tabular}
\end{center}

\index{node@\code{node(id:)}!the singular field alone does it}%
The singular \code{node(id: ID!): Node} in a second subgraph is what silences
the check. \code{nodes}, which differs from it only in returning \code{[Node]}
rather than \code{Node}, does not. Neither does implementing the interface
without declaring the field.

\index{node@\code{node(id:)}!an interface with nothing behind it works too}%
Row four is the one to sit with. \code{Speaker} no longer implements
\code{Node} there, so that subgraph's \code{Node} has no object type behind it
at all and its \code{node} field can return nothing whatsoever. The composer
notices, says so, and then composes the unresolvable graph anyway:

\begin{minted}[fontsize=\footnotesize,breakanywhere]{text}
Subgraph "speakers": The Interface "Node" is used as an output type without at least one Object type implementation
defined in the schema.
\end{minted}

\index{resolvability!a warning about the wrong thing}%
That is a warning about a subgraph with no rows in it, standing in for an error
about three fields that cannot be fetched.

\subsection{Two Fields, Moved Down One File}

\index{resolvability!declaration order decides the answer}%
The order the root fields are written in decides it too, and this is the
demonstration I would show anybody who thinks a passing composition is a
result. Take the Sessions document exactly as the service exports it. Move
\code{node} and \code{nodes} from the top of \code{type Query} to the bottom of
the same type. Same fields, same types, same directives, same document
otherwise. Compose the same three subgraphs again, and here is the path with
the first of its four reason lines:

\begin{minted}[fontsize=\footnotesize,breakanywhere]{text}
The field "averageScore" is unresolvable at the following path:
 query {
  sessions {
   edges {
    node {
     averageScore <--
    }
   }
  }
 }
This is because:
 - The root type field "Query.sessions" is defined in the following subgraph: "sessions".
\end{minted}

\index{resolvability!the control run}%
And with the Ratings key put back the way it belongs, that same reordered
document composes exactly as the original does. So the reordering is not the
defect. The reordering is what decides whether the defect is found.

\index{figure!the order decides}%
Figure~\ref{fig:ch15-the-order-decides} is the two documents side by side.

\begin{figure}[htbp]
  \centering
  % The same four root fields, in two declaration orders, against the same broken
% subgraph. Above, the composer finds nothing. Below, it finds the defect on
% its second step.
%
% Same idiom as figures/tikz/ch01-one-field.tex, which fixes it, and drawn as
% that file's two-lane timeline because the subject here is order and nothing
% else. Direct sibling of figures/tikz/ch07-the-walk.tex: same left-to-right
% axis, same ticks, same cross for a point the walk cannot pass. What is new is
% a third kind of stage, and it needs a mark of its own.
%
% The three marks, and why each is the shape it is:
%
%   a vertical stroke   a field the composer actually walked
%   a cross             the walk reaching a step it cannot take, from ch07
%   a hollow square     a field skipped without being walked
%
% The hollow square is the addition. A skipped stage is not a blocked one and
% must not read as a cross, and it is not a walked one either, so a stroke is
% wrong as well. An empty box says the composer arrived and did nothing, which
% is exactly what the record lookup does. Square rather than round because the
% idiom has no rounded corners anywhere.
%
% The lower lane's axis stops after its second stage rather than running to the
% end. That is not decoration: the composer returns at the first root field it
% cannot walk, so the two fields after it are genuinely never looked at in that
% run, and drawing them would say the opposite.
%
% No quotation marks anywhere in this file, including this comment block, per
% SPEC decision 30 and the note at the head of ch07-the-walk.tex.
%
% Bare tikzpicture (house style): the figure environment, caption and label
% live at the call site in the section file.
\begin{tikzpicture}[
  x=1mm, y=1mm,
  every node/.append style={font=\sffamily\scriptsize},
  trunk/.style={draw, line width=0.4pt},
  axis/.style={draw, line width=0.4pt,
    -{Straight Barb[length=1.4mm, width=1.6mm]}},
  tick/.style={draw, line width=0.4pt},
  skip/.style={draw, line width=0.4pt, minimum width=2.4mm,
    minimum height=2.4mm, inner sep=0pt},
  event/.style={anchor=south, align=center, text width=30mm, inner sep=1pt},
  outcome/.style={anchor=west, align=left, text width=32mm, inner sep=1pt},
  lane/.style={anchor=west, font=\sffamily\scriptsize\itshape},
  spanlabel/.style={anchor=north, align=center, inner sep=2pt},
  span/.style={draw, line width=0.4pt},
]

% ------------------------------------------------------------------- lane A
\node[lane] at (0,14) {A. type Query as the service exports it};

\draw[trunk] (0,0) -- (112,0);

% node and nodes are walked. sessions and sessionById are not.
\foreach \x in {14,44}{\draw[tick] (\x,-1.5) -- (\x,1.5);}
\foreach \x in {74,102}{\node[skip] at (\x,0) {};}

\node[event] at (14,3)  {node\\walked, and passes};
\node[event] at (44,3)  {nodes};
\node[event] at (74,3)  {sessions};
\node[event] at (102,3) {sessionById};

\draw[span] (6,-4) -- (6,-7) -- (52,-7) -- (52,-4);
\node[spanlabel, text width=48mm] at (29,-8)
  {shared with the speakers subgraph. Session is recorded on the way past, and
   the field passes because that subgraph answered for its own type};

\draw[span] (66,-4) -- (66,-7) -- (110,-7) -- (110,-4);
\node[spanlabel, text width=44mm] at (88,-8)
  {Session is already in the record, so neither is walked};

\node[outcome] at (118,0) {composes.\\Nothing printed.};

% ------------------------------------------------------------------- lane B
\begin{scope}[shift={(0,-46)}]
  % Higher than lane A's title: this lane's first event runs to three lines
  % where lane A's runs to two, and at 14 the two collide.
  \node[lane] at (0,18) {B. the same four fields, node and nodes declared last};

  % Stops where the walk stops. Solid to the first stage, dashed into the
  % second, and nothing after it.
  \draw[trunk] (0,0) -- (14,0);
  \draw[axis, dashed] (14,0) -- (44,0);

  \draw[tick] (14,-1.5) -- (14,1.5);
  \draw[tick] (42.8,-1.5) -- (45.2,1.5);
  \draw[tick] (42.8,1.5) -- (45.2,-1.5);

  \node[event] at (14,3) {sessions\\walked first};
  \node[event] at (44,3) {into Session,\\and no further};

  \node[event, text width=30mm, anchor=south] at (74,3)  {node};
  \node[event, text width=30mm, anchor=south] at (102,3) {nodes};
  \node[spanlabel, text width=44mm] at (88,-2)
    {never reached: the check returns at the first root field it cannot walk};

  \draw[span] (6,-6) -- (6,-9) -- (52,-9) -- (52,-6);
  \node[spanlabel, text width=48mm] at (29,-10)
    {nothing has been recorded yet, so the walk goes to ratings and finds no
     key it may enter through};

  \node[outcome] at (118,0) {fails, naming\\Query.sessions.};
\end{scope}

\end{tikzpicture}

  \caption{The same four root fields, in two orders, against the same broken
  subgraph. Above, the two shared fields are checked first: \code{node} records
  \code{Session} on the way past and then passes, because the other subgraph
  declaring it answered for its own type, and the two fields that follow are
  waved through on the strength of that record rather than walked. Below, the
  fields that would have been waved through are checked first, with nothing
  recorded yet, and the first of them fails; the composer returns there, so the
  two fields after it are never looked at. Neither the graph nor the mistake is
  different between the two lanes.}
  \label{fig:ch15-the-order-decides}
\end{figure}

\subsection{Why}

\index{resolvability!the record is written on touch}%
The composer keeps a record of the entities it has walked into, so that a graph
with a cycle in it does not walk forever. An entity in that record is skipped
when the walk reaches it again. What decides whether an entity goes into the
record is being \emph{visited}, and what the code comment beside the check says
is that only a resolvable one can be in there. Those are different claims, and
the gap between them is this whole section.

\index{shareable@\code{"@shareable}!an OR across subgraphs}%
\code{Query.node} is declared by two subgraphs, so the composer treats it as
one field with two answers and accepts it if either of them works. Walking the
Sessions answer records \code{Session} and then fails, because ratings is
unreachable from it. Walking the Speakers answer succeeds, for \code{Speaker},
which has nothing to do with the question. The field passes. \code{Session}
stays in the record, marked as visited and never resolved, and \code{sessions},
\code{sessionById} and the mutation payload all reach it afterwards and are
skipped.

\index{resolvability!vacuous truth}%
Row four of the table above falls out of the same reading. An abstract type
with no implementations has nothing to check, and a walk over nothing succeeds.

\index{resolvability!the reading was tested before it was believed}%
That reading is mine, of a composer somebody else wrote, and it earns its place
here by predicting something. Move the shared field below the unshared ones and
nothing is in the record when they are checked, so the answer should change.
It does, which is the experiment two subsections up. Nobody reorders a schema
for its own sake; the reorder is in this chapter because it is the run that
would have gone the other way if the explanation were wrong.
\index{resolvability!the graph that composes and cannot be served|)}%
